\documentclass[11pt,a4paper]{article}

\usepackage[utf8]{inputenc}
\usepackage[T1]{fontenc}
\usepackage[spanish,es-tabla,es-noshorthands]{babel}
\usepackage{lmodern}
\usepackage{microtype}

\usepackage{amsmath,amssymb,amsfonts}
\usepackage{bm}
\usepackage{graphicx}
\usepackage{booktabs}
\usepackage{array}
\usepackage{multirow}
\usepackage{caption}
\usepackage{subcaption}
\usepackage{float}
\usepackage{enumitem}
\usepackage{xcolor}
\usepackage{listings}
\usepackage[a4paper,margin=2.5cm]{geometry}
\usepackage[hidelinks]{hyperref}
\usepackage{cleveref}

\definecolor{codegray}{rgb}{0.5,0.5,0.5}
\definecolor{codegreen}{rgb}{0.0,0.5,0.0}
\definecolor{codepurple}{rgb}{0.58,0.0,0.82}
\definecolor{backcolour}{rgb}{0.96,0.96,0.96}

\begin{document}
\begin{titlepage}
\centering
{
\includegraphics[width=0.7\textwidth]{Imagenes/logoTEC.png}\par}
\vspace{1cm}
{\bfseries\Large Trabajo Práctico III \par}
\vspace{1cm}
{\scshape\Large Curso:\\
Inteligencia Artificial \par}
\vspace{2cm}
{\scshape\Huge Redes Neuronales \par}
\vspace{2cm}
{\Large Profesor: \\
Saúl Calderón Ramírez\\ \par}
\vfill

\vfill
{\Large Estudiantes \par}
{ \Large Kener Adiel Castillo \\ Alejandro Cerdas Vargas \\ José Pablo Pérez Delgado \par}
\vfill

{\Large Carnet \par}
{ \Large 2023179846 \\ 2023073037 \\ 2023051991\par}

\vfill
{\Large \today \par}
\end{titlepage}

\tableofcontents

\newpage

\section{Introducción}

Las redes neuronales artificiales constituyen un modelo de cómputo inspirado en
la organización del sistema nervioso, capaz de aproximar relaciones no lineales
entre un conjunto de variables de entrada y una variable de salida. El
perceptrón multicapa con una capa oculta es, bajo condiciones suaves, un
aproximador universal, ya que puede aproximar cualquier función continua con error arbitrariamente pequeño, siempre que disponga de
suficientes neuronas ocultas \cite{cybenko1989,hornik1989}. \\

El presente trabajo implementa, desde sus operaciones fundamentales, un
perceptrón de tres capas (entrada, oculta y salida) y lo entrena con el
algoritmo de descenso del gradiente con momentum, una variante
del descenso del gradiente clásico que acelera la convergencia y atenúa las
oscilaciones en direcciones de curvatura elevada. La totalidad del código se
encuentra distribuida en dos archivos:

\begin{itemize}[nosep]
  \item \texttt{Util.py}: utilidades de visualización
        (\texttt{Visualization}), construcción de tablas y exportación a
        \LaTeX{} (\texttt{Table}), y generación de conjuntos de datos
        sintéticos (\texttt{DatasetGenerator}).
  \item \texttt{TP3.ipynb}: la clase \texttt{MultilayerPerceptron} y los
        experimentos de las cuatro partes del enunciado.
\end{itemize}

El documento se organiza siguiendo la estructura del enunciado: la
\Cref{sec:marco} presenta el marco teórico y las ecuaciones implementadas; la
\Cref{sec:parte1} describe la implementación del MLP y sus pruebas unitarias;
las \Cref{sec:parte2,sec:parte3} reportan los experimentos sobre XOR y sobre
datos en $\mathbb{R}^2$; y la \Cref{sec:parte4} aborda la clasificación de
imágenes de fondo de ojo. La \Cref{sec:conclusiones} resume las conclusiones.

\section{Marco teórico} \label{sec:marco}

\subsection{Arquitectura y notación}

Se considera una red de tres capas con $D$ neuronas de entrada, $M$ neuronas en
la capa oculta y $K$ neuronas de salida, descrita en el código por el arreglo
\texttt{neurons\_per\_layer}~$=[D,M,K]$. Sea
$\mathbf{X}\in\mathbb{R}^{N\times D}$ la matriz que agrupa $N$ observaciones de
entrada (una por fila) y $\mathbf{T}\in\mathbb{R}^{N\times K}$ las etiquetas
deseadas. Los parámetros actualizables para mejorar los resultados son:

\begin{align}
  W_o &\in \mathbb{R}^{D \times M}, &
  b_o &\in \mathbb{R}^{1 \times M}
    &&\text{(pesos y sesgos de la capa oculta),}\\
  W_s &\in \mathbb{R}^{M \times K}, &
  b_s &\in \mathbb{R}^{1 \times K}
    &&\text{(pesos y sesgos de la capa de salida).}
\end{align}

\subsection{La función de activación sigmoide}\label{sec:sigmoide}

Entre cada capa lineal se aplica una función de activación no lineal. Si se encadenaran capas puramente lineales, la composición
seguiría siendo una única transformación lineal y la red no podría representar
fronteras de decisión curvas (no podría, por ejemplo, resolver XOR, como se verá
en la sección \ref{sec:parte2}). En este trabajo se emplea la función
sigmoide en ambas capas:
\begin{equation}
  \sigma(z)=\frac{1}{1+e^{-z}},\qquad
  \sigma'(z)=\sigma(z)\,\bigl(1-\sigma(z)\bigr).
  \label{eq:sigmoid}
\end{equation}

La sigmoide tiene tres propiedades que la hacen conveniente aquí. Primero, compacta cualquier número real al intervalo $(0,1)$, de modo que su salida
puede interpretarse directamente como una probabilidad, lo que resulta natural en
clasificación binaria, donde la red estima $P(\text{clase}=1\mid \vec{x})$.
Segundo, es suave y diferenciable en todo punto, requisito para aplicar el
descenso del gradiente. Tercero, su derivada se expresa en términos de la propia
salida, $\sigma'(z)=\sigma(z)(1-\sigma(z))$, lo que permite reutilizar en la
retropropagación los valores ya calculados en la propagación hacia adelante, sin
operaciones adicionales [\ref{eq:delta_s},\ref{eq:delta_o}].

Conviene señalar, por rigor, su principal limitación: como $\sigma'(z)$ alcanza
su máximo $0.25$ y decae a cero en los extremos, en redes muy
profundas el gradiente se atenúa al propagarse por muchas capas saturadas, el llamado desvanecimiento del gradiente, \cite{glorot2010,pascanu2013}.
Por ello las arquitecturas profundas modernas prefieren activaciones como ReLU ($max(0,x)$).
En esta red, de una sola capa oculta, ese inconveniente es no es tan grave, y la
sigmoide se ajusta además a lo solicitado por el enunciado y a su interpretación
probabilística; de ahí su elección.

\subsection{Inicialización de pesos (Xavier/Glorot)}

Antes de poder entrenar la red se puede asignar un valor inicial a los pesos.
Esta elección condiciona si el entrenamiento arranca de forma sana
o si se estanca desde la primera iteración. Conviene entender primero qué
problema resuelve la inicialización de Xavier/Glorot.

Cuando una entrada atraviesa la red, en cada capa se calcula una suma ponderada
de las salidas de la capa anterior [\ref{eq:fwd1}, \ref{eq:fwd2}]. Si los pesos
iniciales son demasiado grandes, la magnitud de esa suma crece capa a
capa y empuja a la sigmoide hacia sus extremos (cercanos a $0$ o a $1$), donde su
derivada $\sigma'(z)=\sigma(z)(1-\sigma(z))$ es casi nula (eq.\ref{eq:sigmoid}); en
consecuencia, durante la retropropagación apenas circula gradiente y la red deja
de aprender. Si, por el contrario, los pesos son
demasiado pequeños, la señal se encoge progresivamente hacia cero y ocurre
el fenómeno simétrico de desvanecimiento. La idea de
Glorot y Bengio \cite{glorot2010} es elegir la escala inicial de los pesos de
modo que la varianza de la señal se conserve aproximadamente constante
al pasar de una capa a la siguiente —tanto hacia adelante (activaciones) como
hacia atrás (gradientes), evitando ambos extremos.

Considérese una neurona que suma $\text{fan\_in}$ entradas independientes (cantidad de entradas que recibe una neurona), cada
una con varianza $\operatorname{Var}(x)$, ponderadas por pesos de media nula y
varianza $\operatorname{Var}(w)$. La varianza de la preactivación resultante es
aproximadamente
\begin{equation}
  \operatorname{Var}(p)\;\approx\;\text{fan\_in}\cdot
     \operatorname{Var}(w)\cdot \operatorname{Var}(x).
  \label{eq:var-prop}
\end{equation}
Para que la varianza de salida iguale a la de entrada en el paso hacia
adelante se necesitaría $\operatorname{Var}(w)=1/\text{fan\_in}$; para preservar
la varianza de los gradientes en el paso hacia atrás se necesitaría, de forma
análoga, $\operatorname{Var}(w)=1/\text{fan\_out}$. Como ambas condiciones no
pueden satisfacerse a la vez (a menos que $\text{fan\_in}=\text{fan\_out}$), se
adopta el promedio armónico de ambas como compromiso:
\begin{equation}
  \operatorname{Var}(w)\;=\;\frac{2}{\text{fan\_in}+\text{fan\_out}}.
  \label{eq:xavier-var}
\end{equation}
Finalmente, para una distribución uniforme $\mathcal{U}(-r,r)$ la varianza vale
$r^2/3$ (resultado de $\frac{(b-a)^2}{12}$). Igualando $r^2/3$ con \cref{eq:xavier-var} y despejando $r$ se obtiene
exactamente la cota empleada en el código:
\begin{equation}
  r_{\ell}=\sqrt{\frac{6}{\text{fan\_in}+\text{fan\_out}}},\qquad
  w \sim \mathcal{U}(-r_\ell,\,r_\ell).
  \label{eq:xavier}
\end{equation}
En el código, la función \texttt{xavier\_glorot\_init} calcula
\cref{eq:xavier}, usándose $r_o=\sqrt{6/(D+M)}$ para $W_o$ y
$r_s=\sqrt{6/(M+K)}$ para $W_s$. \\

Los sesgos no participan del problema de propagación de varianza descrito arriba:
hay un único sesgo por neurona, de modo que no se acumula a lo largo de muchas
conexiones como sí ocurre con los pesos en \cref{eq:var-prop}. Tampoco hace falta
que sean aleatorios para romper la simetría entre neuronas, esa función ya
la cumplen los pesos inicializados al azar, lo que impide que todas las neuronas
ocultas aprendan exactamente lo mismo. Al fijarlos en cero, la preactivación
inicial queda centrada en torno al origen, que es justamente donde la sigmoide
tiene su mayor pendiente ($\sigma'(0)=0.25$); así el gradiente fluye con fuerza
desde la primera iteración. Por estas razones, inicializar los sesgos en cero es
la práctica estándar y la que se adopta aquí.

\subsection{Propagación hacia adelante} \label{sec:forward}

La salida de la red se calcula propagando la entrada a través de las dos capas
(\texttt{forward}):
\begin{align}
  P_o &= X W_o + b_o,
    & Y_o &= \sigma(P_o)
    && \in R^{N \times M},\label{eq:fwd1}\\
  P_s &= Y_o W_s + b_s,
    & Y_s &= \sigma(P_s)
    && \in R^{N \times K}, \label{eq:fwd2}
\end{align}
donde $P$ denota los pesos netos e $Y$ las salidas de cada capa luego de pasar por la activación, en este caso sigmoide. El sesgo se suma sobre las $N$ filas.

\subsection{Función de error} \label{sec:error}

El entrenamiento necesita una medida escalar de cuán lejos está la salida de la
red, $Y_s$, de la etiqueta deseada $T$; esa medida es la
función de error (o pérdida) que el descenso del gradiente minimiza. Se
emplea el error cuadrático medio (Mean Squared Error (MSE)), calculado en
\texttt{get\_error} mediante \texttt{mse\_loss} de PyTorch:
\begin{equation}
  E(Y_s,T)
   = \frac{1}{N K}\sum_{n=1}^{N}\sum_{j=1}^{K}
     \bigl(y_{s,nj}-t_{nj}\bigr)^2 .
  \label{eq:mse}
\end{equation}
Cada término $(y_{s,nj}-t_{nj})^2$ mide la discrepancia entre la salida y el valor
deseado para la muestra $n$ y la neurona de salida $j$. Elevar al cuadrado tiene
dos efectos: hace que el error sea siempre positivo y simétrico, y \emph{penaliza
más las desviaciones grandes} que las pequeñas, lo que orienta el aprendizaje a
corregir primero los errores más groseros; además, $(\cdot)^2$ es suave y
diferenciable, condición necesaria para el descenso del gradiente. \\

El doble promedio no es cosmético. Dividir entre el número de muestras $N$
convierte la suma en un error medio por muestra, de modo que su valor y,
sobre todo, la magnitud del gradiente que de él se deriva, no depende del
tamaño del conjunto ni del lote; sin esa división, procesar más datos
aumentaría artificialmente la pérdida y el paso de actualización, y obligaría a
re-calibrar $\alpha$ cada vez que cambia $N$. Dividir además entre el número de
salidas $K$ lo vuelve un error medio por neurona de salida, lo que hace la
métrica comparable entre la red de $K{=}2$ de XOR (sección \ref{sec:parte2}) y las de
$K{=}1$ de las Partes 3 y 4.

\subsection{Retropropagación de los deltas}\label{sec:backprop}

Para minimizar el error hay que saber cómo cambia $E$ (función de pérdida) al modificar cada peso, es
decir, su gradiente. Calcularlo capa por capa de forma directa sería costoso y
redundante; la retropropagación \cite{rumelhart1986} lo resuelve
aplicando la regla de la cadena de atrás hacia adelante y reutilizando
resultados intermedios. La cantidad central es el delta de cada capa,
$\delta=\frac{\partial E}{\partial P}$, la sensibilidad del error respecto de
la preactivación $P$, mide cuánto contribuye al error cada neurona
antes de la activación, y a partir de él se obtienen los gradientes de
pesos y sesgos.

Para la capa de salida se compone la derivada del error con la de la activación.
Derivando el MSE (\cref{eq:mse}) respecto de $Y_s$ se obtiene un término
proporcional a $(Y_s-T)$, y por la regla de la cadena se
multiplica por la derivada de la sigmoide $\sigma'(P_s)=
Y_s\odot(1-Y_s)$ (\cref{eq:sigmoid}):
\begin{equation}
  \delta_s = (Y_s-T) \odot
                  Y_s \odot (1 - Y_s) \in R^{N \times K}.
  \label{eq:delta_s}
\end{equation}
Para la capa oculta, su preactivación influye en el error a través de la
capa de salida; por eso el delta de salida se retropropaga hacia atrás
multiplicándolo por $W_s^T$ (la misma matriz de la propagación hacia
adelante, transpuesta para invertir el sentido del flujo) y, de nuevo por la regla
de la cadena, por la derivada local de la sigmoide de esta capa:
\begin{equation}
  \delta_o = \bigl(\delta_s W_s^T\bigr)\odot
                  Y_o \odot (1-Y_o) \in R^{N \times M},
  \label{eq:delta_o}
\end{equation}
donde $\odot$ denota el producto elemento a elemento.
El nombre ``retropropagación'' refleja precisamente esta dirección, el error
medido en la salida se propaga hacia las capas anteriores. Estas dos ecuaciones
corresponden exactamente al método \texttt{backpropagate\_deltas}.

\subsection{Actualización con momentum}\label{sec:momentum}

El descenso del gradiente clásico actualiza cada peso dando un paso en la dirección opuesta al gradiente, $W \leftarrow W-\alpha\nabla E$.
Esto tiene un inconveniente, en ``valles'' alargados de la superficie de error (direcciones de mucha curvatura) el gradiente apunta de un lado a otro y el
método oscila, avanzando lento hacia el mínimo. El momentum, que
es lo que solicita este proyecto, mitiga este problema con una analogía física, en lugar de mover el peso directamente con el gradiente, se acumula una velocidad $V$, como la de una bola que rueda por la superficie y gana inercia \cite{polyak1964,qian1999}. La velocidad es una media móvil exponencial de los gradientes, las componentes consistentes (que apuntan
época tras época en la misma dirección) se refuerzan y aceleran la marcha,
mientras que las componentes oscilatorias se cancelan entre sí. El
coeficiente $\gamma\in[0,1]$ regula cuánta inercia se conserva
(\texttt{gamma}). En la formulación implementada (\texttt{update\_weights}) la
velocidad combina la inercia y el gradiente promediado sobre las $N$
muestras:
\begin{align}
  V_{W_o} &\leftarrow \gamma\,V_{W_o}
     + \alpha\,(1-\gamma)\,\frac{1}{N}\,X^T,\delta_o, &
  W_o &\leftarrow \operatorname{clamp}\!\bigl(W_o-V_{W_o}\bigr),
  \label{eq:upd_wo}\\
  V_{W_s} &\leftarrow \gamma\,V_{W_s}
     + \alpha\,(1-\gamma)\,\frac{1}{N}\,Y_o^T,\delta_s, &
  W_s &\leftarrow \operatorname{clamp}\!\bigl(W_s-V_{W_s}\bigr),
  \label{eq:upd_ws}
\end{align}
y de forma análoga para los sesgos, cuyo gradiente es la suma de los deltas
sobre el batch:
\begin{equation}
  V_{b}\leftarrow \gamma\,V_{b}
   + \alpha\,(1-\gamma)\,\frac{1}{N}\sum_{n=1}^{N}\delta_{n},\qquad
  b\leftarrow \operatorname{clamp}(b-V_{b}).
  \label{eq:upd_bias}
\end{equation} \\

El operador $\operatorname{clamp}$ fuerza a cada peso a permanecer en el rango
$[-w_{\max},w_{\max}]$, donde $w_{\max}$ es el parámetro \texttt{max\_weights}.
Actúa como un acotador, si por una tasa de aprendizaje
demasiado alta, un momentum elevado o un gradiente puntualmente grande la
actualización empujara los pesos a crecer sin control, estos podrían
``explotar'' hacia magnitudes enormes (fenómeno del gradiente
explosivo \cite{pascanu2013}), saturando las sigmoides y produciendo valores no
representables (\texttt{inf}/\texttt{NaN}) que arruinarían el entrenamiento. El
recorte garantiza que, aun en el peor caso, los pesos queden acotados y el
proceso no diverja. \\

Conviene destacar dos modificaciones:
\begin{enumerate}
  \item El factor $(1-\gamma)$ mantiene acotado el tamaño efectivo del paso, el
        paso asintótico de una dirección constante es
        $\frac{\alpha(1-\gamma)}{1-\gamma}=\alpha$, por lo que $\alpha$ conserva
        su interpretación de tasa de aprendizaje aun variando $\gamma$
  \item El gradiente se promedia (no se suma) sobre las $N$ muestras
        mediante el factor $\frac{1}{N}$. Esta elección es deliberada y
        consistente con la normalización del MSE (\Cref{sec:error}), sin ella, la
        magnitud del gradiente, y por tanto el paso real, escalaría con el número
        de muestras, de modo que el mismo $\alpha$ se comportaría de forma muy
        distinta con $4$ observaciones (XOR) que con miles (datos en
        $\mathbb{R}^2$ o lotes de imágenes). Al promediar, la dinámica se vuelve independiente del tamaño del lote y $\alpha$ es transferible entre
        experimentos; en particular, valores moderados como el $\alpha=0.01$.
\end{enumerate}

Cuando $\gamma=0$ las ecuaciones de eq.\ref{eq:upd_wo} a eq.\ref{eq:upd_bias} se
reducen al descenso del gradiente estándar
$W \leftarrow W-\alpha\,\nabla_{\!W}E$.

\section{Parte 1: Implementación del perceptrón de tres capas}\label{sec:parte1}

\subsection{Pruebas unitarias}\label{sec:pruebas1}

La verificación del funcionamiento correcto de la red se realiza de forma ordenada, para ir cumpliendo dependencias entre funciones. 

\subsubsection{Función \texttt{forward()}} 
\paragraph{Objetivo:} Verificar que el método forward() ejecuta correctamente la propagación hacia adelante para distintas dimensiones entre las capas de entrada, oculta y salida, produciendo activaciones y tamaños de tensor acordes con la arquitectura de la red.

\paragraph{Entrada:} Para esta prueba unitaria se eligío una red neuronal de $3$, $4$, $1$ neuronas en la capa de entrada, en la capa oculta y en la capa de salida, respectivamente. El conjunto de datos definido es
\begin{align}
  X &= 
  \begin{bmatrix} 
      1.0, 2.0, 3.0 \\
      5.0, 6.0, 7.0 \\
      8.0, 9.0, 10.0
  \end{bmatrix} \\
  T &= 
  \begin{bmatrix}
    0 \\
    1 \\
    0
  \end{bmatrix}  
\end{align}
A su vez como la función forward depende de $W_o \,, W_s \,, b_o \,, b_s$, por lo tanto para esta prueba también se define
\begin{align}
  W_o &= 
  \begin{bmatrix}
    0.1, 0.2, 0.3, 0.4 \\
    0.5, 0.6, 0.7, 0.8 \\
    0.9, 1.0, 1.1, 1.2
  \end{bmatrix} \\
  W_s &=
  \begin{bmatrix}
    2.0 \\
    2.1 \\
    2.2 \\
    2.3    
  \end{bmatrix} \\
  b_o &= 
  \begin{bmatrix}
    3.0, 3.1, 3.2, 3.3
  \end{bmatrix} \\
  b_s &= 
  \begin{bmatrix}
    0.5
  \end{bmatrix}
\end{align} 

\paragraph{Salidas esperadas:}

Como se muestra en la sección \ref{sec:forward}, se debe realizar las operaciones intermedias que lleven a la salida esperada
\begin{align}
  P_o &= X W_o + B_o \\
  P_o &=
  \begin{bmatrix}
    6.8,  7.5,  8.2,  8.9 \\
    12.8, 14.7, 16.6, 18.5 \\
    17.3, 20.1, 22.9, 25.7
  \end{bmatrix} 
\end{align}
\begin{align}  
  Y_o &= \sigma(P_o) \\
  Y_o &=
  \begin{bmatrix}
    0.9989, 0.9994, 0.9997, 0.9999 \\
    1, 1, 1, 1 \\
    1, 1, 1, 1
  \end{bmatrix} \label{eq:test_Y_o}  
\end{align}
\begin{align}
  P_s &= Y_o W_s + b_s \\
  P_s &= 
  \begin{bmatrix}
    9.0957 \\
    9.1 \\
    9.1
  \end{bmatrix}
\end{align}
\begin{align}
  Y_s &= \sigma(P_s) \\
  Y_s &= 
  \begin{bmatrix}
    0.9999 \\
    0.9999 \\
    0.9999
  \end{bmatrix} \label{eq:esperado_1}
\end{align}  
Por lo tanto el resultado esperado es lo que se ve en \cref{eq:esperado_1}, que representa la predicción.

\paragraph{Resultados:} La prueba se supera si la diferencia absoluta entre los elementos correspondientes de la salida esperadada y el resultado, sea menor a $10^{-3}$. El resultado obtenido es
\begin{equation}
  R = 
  \begin{bmatrix}
    0.9999 \\
    0.9999 \\
    0.9999
  \end{bmatrix}
\end{equation}
Por lo tanto la prueba se considera exitosa.

\subsubsection{Función \texttt{backpropagate\_deltas()}}

\paragraph{Objetivo:} Verificar que el método \texttt{backpropagate\_deltas()} implementa correctamente el algoritmo de retropropagación, calculando los gradientes de las capas de salida y oculta a partir de los valores obtenidos del \texttt{forward()}, manteniendo la consistencia dimensional para los diferentes números de neuronas entre las capas de la red neuronal.

\paragraph{Entrada:} Para esta prueba unitaria se eligío una red neuronal de $3$, $4$, $1$ neuronas en la capa de entrada, en la capa oculta y en la capa de salida, respectivamente. El conjunto de datos definido es
\begin{align}
  X &= 
  \begin{bmatrix} 
      1.0, 2.0, 3.0 \\
      5.0, 6.0, 7.0 \\
      8.0, 9.0, 10.0
  \end{bmatrix} \\
  T &= 
  \begin{bmatrix}
    0 \\
    1 \\
    0
  \end{bmatrix}  
\end{align}

Como se muestra en \cref{eq:delta_s}  y \cref{eq:delta_o}, es necesario definir $Y_s \,, W_s \,, Y_o$, por lo se reutilizarn los cálculos hechos en \cref{eq:test_Y_o} y \cref{eq:esperado_1}, entonces
\begin{align}
  W_s &=
  \begin{bmatrix}
    2.0 \\
    2.1 \\
    2.2 \\
    2.3    
  \end{bmatrix} \\
  Y_o &=
  \begin{bmatrix}
    0.9989, 0.9994, 0.9997, 0.9999 \\
    1, 1, 1, 1 \\
    1, 1, 1, 1
  \end{bmatrix} \\
  Y_s &= 
  \begin{bmatrix}
    0.9999 \\
    0.9999 \\
    0.9999
  \end{bmatrix}
\end{align}

\paragraph{Salidas esperadas:} Para obtener la salida esperada $\delta_s$ se debe seguir \cref{eq:delta_s}
\begin{align}
  \delta_s &= (Y_s - T) \odot Y_s \odot (1 - Y_s) \\
  \delta_s &= 
  \begin{align}
        
  \end{align} 
\end{align}

\subsection{Variantes por lotes}

Para los experimentos con imágenes (\Cref{sec:parte4}), donde el conjunto no
cabe en una sola pasada o es muy costoso realizarlo de esta forma, se añadieron los métodos \texttt{train\_batches} y
\texttt{evaluate\_batches}, que iteran sobre un \texttt{DataLoader} de PyTorch
acumulando el error medio por época. La función \texttt{get\_accuracy} distingue
el caso binario ($K=1$, umbralizando en $0.5$) del caso multiclase ($K>1$,
mediante $\arg\max$).

\section{Parte 2: Clasificación de la función XOR}\label{sec:parte2}

\subsection{Planteamiento del problema}

La función lógica XOR sobre entradas booleanas
$\vec{x}\in\{0,1\}^2$ es el ejemplo canónico de un problema no separable
linealmente, es decir, no existe un hiperplano que separe las clases, por lo que un
perceptrón simple no puede resolverlo. La incorporación de una capa oculta
permite construir fronteras de decisión no lineales y resolver la tarea.

\subsection{Arquitectura propuesta}

Se emplea una red
$$D{=}2, M{=}2, K{=}2$$
con codificación one-hot de las dos clases en la salida:
$$T={[1,0],[0,1],[0,1],[1,0]}$$
para las entradas
$$X={[0,0],[0,1],[1,0],[1,1]}$$
respectivamente (clase 0 $=$ XOR falso y clase 1
$=$ XOR verdadero). La figura \ref{fig:xor-img} ilustra la topología.

\begin{figure}[H]
  \centering
  \includegraphics[width=0.8\linewidth]{Imagenes/xor/xor-arg.png}
  \caption{Arquitectura utilizada para resolver el problema de XOR con una red neuronal o perceptrón multicapa, con 2 neuronas de capa de entrada, 2 neuronas de capa oculta y 2 neuronas de capa de salida}
  \label{fig:xor-img}
\end{figure}

\subsection{Búsqueda de hiperparámetros}

Los hiperparámetros $\alpha$ y $\gamma$ se calibran con
Optuna \cite{akiba2019optuna} (función \texttt{objective\_xor}), minimizando el
error final tras un máximo de $20000$ épocas sobre las cuatro observaciones,
con paro temprano en cuanto el error cae por debajo de $\varepsilon$ (los
ensayos exitosos terminan, en la práctica, mucho antes del límite). El espacio
de búsqueda es $\alpha\in[1,200]$ (escala logarítmica) y
$\gamma\in[0,0.95]$, con $w_{\max}=1000$ y $50$ ensayos; la justificación de
estos rangos se desarrolla en la \Cref{sec:convergencia}. Para cada ensayo que
alcanza error $<\varepsilon=10^{-5}$ se registran el número de iteraciones hasta
exactitud $>0.99$ y hasta error $<\varepsilon$, y se grafican la frontera de
decisión y la curva de error (\Cref{fig:xor-result}). La
\Cref{tab:xor-best} resume las configuraciones que lograron error nulo sobre las
cuatro observaciones.

\begin{table}[H]
\centering
\caption{Mejores configuraciones para XOR (error $<\varepsilon$ sobre las
cuatro observaciones). Generada por \texttt{show\_best\_trial} en
\texttt{Optimizacion.tex}.}
\label{tab:xor-best}
\begin{tabular}{@{}cccccc@{}}
\toprule
Intento & $\alpha$ & $\gamma$ & Valor (error) & Iter. precisión & Iter. error\\
\midrule
\multicolumn{6}{c}{\itshape (valores producidos por la corrida de Optuna)}\\
\bottomrule
\end{tabular}
\end{table}

\begin{figure}[H]
  \centering
  \begin{subfigure}{0.48\linewidth}
    \centering
    \IfFileExists{figuras/xor_frontera.png}
      {\includegraphics[width=\linewidth]{figuras/xor_frontera.png}}
      {\fbox{\parbox[c][3.2cm][c]{0.9\linewidth}{\centering\itshape
        \texttt{figuras/xor\_frontera.png}}}}
    \caption{Regiones de decisión aprendidas.}
  \end{subfigure}\hfill
  \begin{subfigure}{0.48\linewidth}
    \centering
    \IfFileExists{figuras/xor_error.png}
      {\includegraphics[width=\linewidth]{figuras/xor_error.png}}
      {\fbox{\parbox[c][3.2cm][c]{0.9\linewidth}{\centering\itshape
        \texttt{figuras/xor\_error.png}}}}
    \caption{Curva de error vs. iteraciones.}
  \end{subfigure}
  \caption{Resultado del entrenamiento sobre XOR para la mejor configuración
    hallada (\textbf{VALORES DE LA MEJOR CORRIDA}). A la izquierda, la frontera no lineal separa correctamente las dos
    clases; a la derecha, el error MSE decae hasta valores inferiores a
    $\varepsilon$.}
  \label{fig:xor-result}
\end{figure}

\subsection{Convergencia}\label{sec:convergencia}

\paragraph{Criterio de paro: error objetivo, no estancamiento.}
Con salida sigmoide el error nunca llega exactamente a cero: $\sigma(z)\in(0,1)$
es asintótica, de modo que siempre existe un residuo $|y-t|>0$. La exactitud del
$100\%$, en cambio, se alcanza muy temprano (basta que cada salida quede del
lado correcto de $0.5$). Por ello el entrenamiento se detiene cuando el error
cae por debajo del objetivo $\varepsilon$ (paro temprano por objetivo,
parámetro \texttt{epsilon} de \texttt{train\_mlp}), lo que además abarata la
búsqueda: los ensayos exitosos no consumen el presupuesto completo de épocas.
Conviene advertir que el criterio alternativo de parar por
\emph{estancamiento} (cambio de error entre épocas menor que un umbral) es
inadecuado en este problema: la superficie de error de XOR presenta un
\emph{plateau} inicial con $E\approx 0.25$ (todas las salidas $\approx 0.5$) en
el que el error varía de forma minúscula por época sin que la red haya
aprendido nada; un paro por estancamiento aborta ahí el entrenamiento y sesga
la calibración hacia valores de $\alpha$ artificialmente enormes, los únicos
capaces de atravesar el plateau antes de activar el criterio.

\paragraph{Por qué $\varepsilon=10^{-5}$ es un objetivo exigente.}
Clasificar XOR y alcanzar $\text{MSE}<10^{-5}$ son metas de dificultad muy
distinta. Lo segundo exige $|y-t|\lesssim\sqrt{\varepsilon}\approx 3\times
10^{-3}$ en cada salida, es decir, preactivaciones $|P_s|\gtrsim 6$ y pesos un
orden de magnitud mayores que los iniciales de Xavier (\cref{eq:xavier}).
Además, el gradiente de salida (\cref{eq:delta_s}) contiene el factor
$Y_s\odot(1-Y_s)$, que se desvanece justamente al saturar: si el residuo es
$\epsilon_r=|y-t|$, el delta escala como $\epsilon_r^2$, por lo que cada orden
de magnitud adicional de precisión cuesta aproximadamente $10$ veces más
épocas. Empíricamente (con $\alpha=20$, $\gamma=0.9$): exactitud $100\%$ en
$\sim 10^2$ épocas, $\text{MSE}<10^{-3}$ en $\sim 3\times 10^2$,
$\text{MSE}<10^{-4}$ en $\sim 2.5\times 10^3$ y $\text{MSE}<10^{-5}$ recién en
$\sim 2.3\times 10^4$. En este régimen el producto
$\alpha\times\text{épocas}$ es aproximadamente constante
($\approx 4.7\times 10^5$), lo que fija un piso práctico para $\alpha$: con el
presupuesto de $2\times 10^4$ épocas se requiere $\alpha\gtrsim 25$. Este
``$\alpha$ grande'' no es síntoma de error, sino consecuencia de combinar MSE
con sigmoide (con entropía cruzada el factor $Y_s(1-Y_s)$ se cancela y
bastarían tasas de aprendizaje un orden de magnitud menores).

\paragraph{Elección del rango de $\alpha$.}
No existe un rango de $\alpha$ ``normal'' universal: la tasa de aprendizaje
solo tiene significado relativo a la escala del gradiente, que aquí depende de
la pérdida (MSE encoge el gradiente frente a entropía cruzada), del promedio
$1/N$ (\cref{eq:upd_wo}) y de la forma del momentum. El rango se eligió con dos
criterios prácticos: (i) un barrido logarítmico corto de pocas épocas para
identificar el mayor $\alpha$ que aún desciende de forma estable, y (ii)
verificar \emph{a posteriori} que los mejores ensayos de Optuna caigan en el
interior del rango y no acumulados contra un borde —un óptimo pegado al borde
indica que el verdadero óptimo está fuera del rango (o que el objetivo premia
valores extremos por un defecto del criterio de paro, como ocurría con el paro
por estancamiento). Con $\alpha\in[1,200]$ los ensayos exitosos se concentran
en $\alpha\in[29,96]$, cómodamente en el interior.

\paragraph{El factor $(1-\gamma)$: ventajas y riesgos.}
La formulación con promedio móvil (\cref{eq:upd_wo}) hace que el paso
estacionario sea $\approx\alpha\,\nabla E$ \emph{independientemente} de
$\gamma$, desacoplando ambos hiperparámetros: $\alpha$ controla el tamaño del
paso y $\gamma$ solo la inercia, lo que es deseable cuando Optuna los busca
simultáneamente (en el momentum clásico $V\leftarrow\gamma V+\alpha\nabla E$ el
paso estacionario es $\alpha/(1-\gamma)$, de modo que cambiar $\gamma$ cambia
la tasa efectiva). Tiene, sin embargo, dos contrapartidas. Primero, el
transitorio es más lento: los primeros pasos quedan escalados por $(1-\gamma)$
hasta que la velocidad acumula. Segundo, el caso límite $\gamma\to 1$ anula por
completo la actualización ($\alpha(1-\gamma)\to 0$), por lo que el espacio de
búsqueda se acota a $\gamma\le 0.95$. Cabe notar que \texttt{optim.SGD} de
PyTorch (usado para AlexNet en la \Cref{sec:parte4}) implementa la forma
clásica, así que los valores de $\alpha$ no son directamente comparables entre
ambos modelos.

\paragraph{Elección de $w_{\max}$.}
El recorte de pesos solo debe actuar como red de seguridad ante pasos
explosivos, nunca como restricción activa: si $w_{\max}$ fuera menor que los
pesos que la solución requiere ($|w|\approx 13$ como máximo en las corridas
exitosas), el objetivo $\varepsilon$ sería directamente inalcanzable, y un peso
``pegado'' a $\pm w_{\max}$ tras el entrenamiento delataría que el recorte está
distorsionando el aprendizaje. Con $w_{\max}=1000$ el recorte no se activa en
ninguna corrida sana y solo intervendría ante una divergencia numérica.

\paragraph{Lectura de los resultados.}
El atributo \texttt{trial.user\_attr["error"]} registra la primera iteración en
que el error cae por debajo de $\varepsilon$, lo que documenta la cantidad de
iteraciones necesarias para la convergencia de cada configuración exitosa. En
la \Cref{tab:xor-best} se observa que las configuraciones con $\gamma>0$
(momentum activo) convergen en menos iteraciones que aquellas con
$\gamma\approx 0$, en concordancia con el efecto acelerador del término de
inercia descrito en \Cref{sec:momentum}; asimismo, dentro del rango explorado,
mayor $\alpha$ reduce las iteraciones necesarias, consistente con la relación
$\alpha\times\text{épocas}\approx\text{constante}$ del régimen saturado.

%==============================================================================
\section{Parte 3: Clasificador de datos en \texorpdfstring{$\mathbb{R}^2$}{R2}}\label{sec:parte3}
%==============================================================================

\subsection{Generación de los conjuntos de datos}

La clase \texttt{DatasetGenerator} de \texttt{Util.py} produce dos conjuntos
sintéticos de $N=10{,}000$ muestras mediante \textsc{scikit-learn}:

\begin{itemize}[nosep]
  \item \textbf{Linealmente separable} (\texttt{generate\_linearly\_separable}):
        dos cúmulos gaussianos con \texttt{make\_blobs}, dos centros y
        desviación estándar $0.35$.
  \item \textbf{No linealmente separable}
        (\texttt{generate\_nonlinearly\_separable}): dos medias lunas
        entrelazadas con \texttt{make\_moons} y ruido $0.08$.
\end{itemize}

La \Cref{fig:r2-datos} muestra ambos conjuntos. La partición en entrenamiento
($80\%$) y validación ($20\%$) se realiza con \texttt{train\_test\_split}
(\textsc{scikit-learn}), con barajado y semilla controlada para reproducibilidad
(\texttt{DatasetGenerator.split\_data}).

\begin{figure}[htbp]
  \centering
  \begin{subfigure}{0.48\linewidth}
    \centering
    \IfFileExists{figuras/r2_separable.png}
      {\includegraphics[width=\linewidth]{figuras/r2_separable.png}}
      {\fbox{\parbox[c][3.2cm][c]{0.9\linewidth}{\centering\itshape
        \texttt{figuras/r2\_separable.png}}}}
    \caption{Conjunto linealmente separable (\texttt{make\_blobs}).}
  \end{subfigure}\hfill
  \begin{subfigure}{0.48\linewidth}
    \centering
    \IfFileExists{figuras/r2_noseparable.png}
      {\includegraphics[width=\linewidth]{figuras/r2_noseparable.png}}
      {\fbox{\parbox[c][3.2cm][c]{0.9\linewidth}{\centering\itshape
        \texttt{figuras/r2\_noseparable.png}}}}
    \caption{Conjunto no separable (\texttt{make\_moons}).}
  \end{subfigure}
  \caption{Conjuntos de datos sintéticos en $\mathbb{R}^2$ usados en la
    \Cref{sec:parte3}.}
  \label{fig:r2-datos}
\end{figure}

\subsection{Configuraciones evaluadas}

Se entrena una red $D{=}2\to M\to K{=}1$ con $\alpha=0.01$ y $w_{\max}=10^5$,
variando el número de neuronas ocultas $M\in\{2,20\}$ y el coeficiente de
\emph{momentum} $\gamma\in\{0,\,0.9\}$, lo que da las cuatro configuraciones de
la clase \texttt{Analyze}:
\begin{equation*}
  (M,\gamma)\in\bigl\{(20,0.9),\;(20,0),\;(2,0.9),\;(2,0)\bigr\}.
\end{equation*}

\paragraph{Interpretación de $P=50$.}
El enunciado solicita ``$P=50$ iteraciones''. La Parte~4 desambigua este símbolo
al referirse a ``la cantidad de iteraciones $P$'' y fijar luego ``$P=50$'' para
la calibración con Optuna, por lo que $P$ designa el número de épocas/repeticiones
del experimento. Para obtener curvas de aprendizaje robustas y comparables se
adopta la interpretación de \textbf{$P=50$ repeticiones} de la partición
$80/20$: cada configuración se entrena $50$ veces con semillas distintas y se
promedia la curva iteraciones-vs-error, tal como implementa
\texttt{run\_50\_iterations} (y la variante vectorizada
\texttt{run\_average\_curves}). Esta decisión se documenta explícitamente porque
satisface simultáneamente las tres exigencias del enunciado: eje $x$ en
iteraciones, $P=50$ y partición $80/20$ con \textsc{scikit-learn}.

\subsection{Curvas de aprendizaje}

La \Cref{fig:r2-curvas} presenta la tabla de gráficas iteraciones-vs-error (una
por configuración) solicitada en el apartado~3.a$'$. Las curvas de
entrenamiento y validación se muestran superpuestas para detectar
sobreajuste.

\safefig{figuras/r2_curvas_2x2.png}{0.95\linewidth}{%
  Tabla de curvas de aprendizaje (error MSE vs. iteraciones) para las cuatro
  configuraciones $(M,\gamma)$, promediadas sobre $P=50$ particiones $80/20$.
  Cada panel superpone la curva de entrenamiento y la de validación.}{fig:r2-curvas}

\subsection{Error final de convergencia}

La \Cref{tab:r2-final} reporta el error final de convergencia por configuración,
generada por la rutina de análisis y exportada a
\texttt{Comparación\_Configuraciones.tex}.

\begin{table}[H]
\centering
\caption{Error MSE final de convergencia por configuración (apartado 3.b$'$).
Datos exportados a \texttt{Comparación\_Configuraciones.tex}.}
\label{tab:r2-final}
\begin{tabular}{@{}lcc@{}}
\toprule
Configuración & Error (validación) & Error (entrenamiento)\\
\midrule
$M=20,\ \gamma=0.9$ & $0.252442$ & $0.252159$\\
$M=20,\ \gamma=0$   & $0.239532$ & $0.239002$\\
$M=2,\ \gamma=0.9$  & $0.251695$ & $0.251072$\\
$M=2,\ \gamma=0$    & $0.270454$ & $0.271777$\\
\bottomrule
\end{tabular}
\end{table}

\subsection{Discusión: relación entre \texorpdfstring{$\alpha$}{alpha},
\texorpdfstring{$M$}{M} y el error}\label{sec:disc-r2}

Esta sección responde a las preguntas del apartado~3.b, integrando la evidencia
empírica con la literatura.

\paragraph{(1) Ajuste de $\alpha$ en función de $M$.}
Al aumentar $M$ crece la cantidad de parámetros (los pesos $\mathbf{W}_o$ son
$D\times M$ y $\mathbf{W}_s$ son $M\times K$) y, por tanto, la magnitud del
gradiente acumulado en la capa de salida, cuya entrada neta es una suma de $\sim
M$ términos. Si la varianza de pesos y activaciones es comparable, esa suma
escala aproximadamente con $\sqrt{M}$; para que el tamaño efectivo del paso no
crezca con la red conviene \textbf{reducir $\alpha$ a medida que $M$ aumenta}
(del orden de $1/\sqrt{M}$ como regla gruesa). Esta es la misma motivación
detrás de la inicialización Xavier/Glorot~\cite{glorot2010} de
\cref{eq:xavier}. Además, el \emph{momentum} amplifica el paso efectivo en
direcciones consistentes en un factor $\approx 1/(1-\gamma)$, por lo que con
$\gamma=0.9$ conviene un $\alpha$ aún menor que con $\gamma=0$ para una misma
$M$.

\paragraph{(2) Relación entre $M$ y el error de entrenamiento.}
Por el teorema de aproximación universal~\cite{cybenko1989,hornik1989}, una red
de una capa oculta con suficientes neuronas puede aproximar cualquier función
continua; en consecuencia, \emph{a mayor $M$, menor error de entrenamiento
alcanzable}. Un $M$ demasiado pequeño produce \emph{subajuste} (el error se
estanca alto por falta de capacidad), mientras que un $M$ excesivo reduce el
error de entrenamiento pero puede inducir \emph{sobreajuste}, aumentando el
error de validación. Este es el clásico compromiso sesgo-varianza~\cite{geman1992}:
$M$ debe elegirse buscando el mínimo del error de \emph{validación}, no del de
entrenamiento. En la \Cref{tab:r2-final} se observa que las cuatro
configuraciones convergen a errores próximos ($\approx 0.24$--$0.27$): con
$\alpha=0.01$ y un presupuesto de iteraciones acotado, la diferencia de
capacidad entre $M=2$ y $M=20$ no llega a manifestarse plenamente, lo que es
coherente con el punto~(1) —un $\alpha$ pequeño limita el avance, especialmente
para $M$ grande.

\paragraph{(3) Elección de $M$ en función de $D$.}
No existe una fórmula óptima cerrada, pero la literatura ofrece reglas
heurísticas para fijar un punto de partida de $M$ a partir de $D$ y $K$:
\begin{itemize}[nosep]
  \item \textbf{Pirámide geométrica}~\cite{masters1993}:
        $M\approx\operatorname{round}\!\bigl(\sqrt{D\,K}\bigr)$.
  \item \textbf{Cota de Kolmogorov/Hecht-Nielsen}~\cite{hechtnielsen1987}:
        $M\le 2D+1$.
  \item \textbf{Reglas prácticas}~\cite{heaton2008}: $K\le M\le 2D$, o bien
        $M=\tfrac{2}{3}D+K$.
  \item \textbf{Cotas según el número de muestras}~\cite{huang1998}: una red de
        una capa oculta con activación acotada aprende $N$ patrones con a lo
        sumo $N$ neuronas; en la práctica $M\ll N$ para evitar sobreajuste.
\end{itemize}
Las revisiones de Stathakis~\cite{stathakis2009} y Karsoliya~\cite{karsoliya2012}
coinciden en que estas reglas sólo proveen un punto de partida que debe
refinarse empíricamente (por validación o con Optuna, como en este trabajo). En
la \Cref{sec:parte4}, con imágenes aplanadas, $D=d^2\cdot 3$ es muy grande, por
lo que $\sqrt{D\,K}$ ya sugiere valores de $M$ elevados; conviene además acotar
$M$ por costo computacional y por la escasez de muestras ($\sim 564$ de
entrenamiento) para no sobreajustar.

%==============================================================================
\section{Parte 4: Clasificador de imágenes de fondo de ojo (ACRIMA)}\label{sec:parte4}
%==============================================================================

\subsection{Conjunto de datos}

Se emplea el conjunto ACRIMA~\cite{diazpinto2019acrima} de imágenes de fondo de
ojo etiquetadas como \emph{Glaucoma} ($396$ imágenes) y \emph{Non Glaucoma}
($309$ imágenes), cargado con \texttt{ImageFolder} de
\texttt{torchvision}. Las imágenes se redimensionan y se particionan en
entrenamiento ($80\%$) y prueba ($20\%$) con \texttt{random\_split} y semilla
fija (clase \texttt{Dataset}).

\subsection{Normalización \texorpdfstring{$\ell_\infty$}{l-infinito}}\label{sec:parte4-norm}

Cada imagen se normaliza al rango $[0,1]$ dividiendo por su valor máximo
(normalización por observación, norma $\ell_\infty$):
\begin{equation}
  \mathbf{X}_{\text{norm}}=\frac{\mathbf{X}}{\max(\mathbf{X})+\epsilon},
  \qquad \epsilon=10^{-10},
  \label{eq:linf}
\end{equation}
implementada en \texttt{ImageNormalizerMLP} (que aplana la imagen a un vector de
$D=d^2\cdot 3$ componentes) y \texttt{ImageNormalizerAlex} (que reordena los
ejes a formato $C\times H\times W$ para AlexNet). El término $\epsilon$ evita la
división por cero.

\paragraph{Pruebas unitarias de la normalización.}
Se documentan dos pruebas (\Cref{tab:pruebas}):
\begin{itemize}[nosep]
  \item \texttt{test\_normalizer\_rango\_0\_1}: verifica que la salida queda
        contenida en $[0,1]$ para una imagen aleatoria.
  \item \texttt{test\_normalizer\_max\_es\_1}: coloca un único píxel con el
        valor máximo y verifica que, tras dividir por el máximo, ese píxel pasa
        a valer exactamente $1$.
\end{itemize}

\subsection{Calibración del MLP}

Se definen tres configuraciones de resolución de entrada,
$d\in\{32,48,64\}$, que dan $D=d^2\cdot 3\in\{3072,\,6912,\,12288\}$ neuronas de
entrada y $K=1$ de salida. Para cada una, \texttt{study\_acrima} calibra
$\alpha$, $\gamma$ y $M$ con Optuna sobre $P=50$ iteraciones, con
$M\in[d,\,570]$ (el límite superior corresponde al total de imágenes,
consistente con la cota~\cite{huang1998} discutida en \Cref{sec:disc-r2}).
La \Cref{fig:acrima-mlp} muestra la curva de error del MLP con los mejores
hiperparámetros hallados.

\safefig{figuras/acrima_mlp_error.png}{0.7\linewidth}{%
  Curva de error de entrenamiento del MLP sobre ACRIMA con los hiperparámetros
  calibrados por Optuna ($\alpha\!\approx\!1.561$, $\gamma\!\approx\!0.626$,
  $M\!=\!118$, $d\!=\!64$).}{fig:acrima-mlp}

\subsection{Arquitectura AlexNet y su adaptación}

Se reutiliza AlexNet~\cite{krizhevsky2012alexnet} con pesos preentrenados en
ImageNet (\texttt{models.alexnet(weights=AlexNet\_Weights.DEFAULT)}),
aplicando \emph{transferencia de aprendizaje}. La \textbf{modificación necesaria}
a la arquitectura original consiste en reemplazar la última capa del
clasificador, originalmente \texttt{Linear(4096,\,1000)} para las $1000$ clases
de ImageNet, por una capa \texttt{Linear(4096,\,1)} adecuada a la clasificación
binaria, seguida de una activación sigmoide (\cref{eq:sigmoid}) para producir una
probabilidad. La entrada se mantiene en $224\times 224\times 3$, normalizada
según \cref{eq:linf}. El entrenamiento usa el optimizador SGD con \emph{momentum}
(\texttt{optim.SGD}, $\text{lr}=\alpha$, $\text{momentum}=\gamma$) y pérdida MSE
(\texttt{nn.MSELoss}), calibrando $\alpha$ y $\gamma$ con Optuna
(\texttt{study\_alex\_acrima}).

\begin{lstlisting}[caption={Adaptación de la cabeza de AlexNet para clasificación binaria.},label=lst:alex]
model = models.alexnet(weights=models.AlexNet_Weights.DEFAULT)
num_ftrs = model.classifier[6].in_features      # 4096
model.classifier[6] = nn.Linear(num_ftrs, 1)    # 1000 -> 1 (binaria)
# entrenamiento: salida = Sigmoid(model(x)),  perdida = MSE,  SGD(lr=alpha, momentum=gamma)
\end{lstlisting}

\subsection{Comparación MLP vs. AlexNet}

Siguiendo el apartado~4.2.b, se realizan $10$ particiones $80/20$ y se evalúa,
para cada una, el mejor MLP completamente conectado frente a la arquitectura
convolucional (\texttt{run\_many}). La \Cref{tab:comp-final} reproduce los
resultados exportados a \texttt{Comparacion\_Final.tex}; cada celda indica
\emph{(error, exactitud)} sobre el conjunto de prueba.

\begin{table}[H]
\centering
\caption{Comparación MLP vs. AlexNet sobre $10$ particiones $80/20$ de ACRIMA.
Cada celda reporta (error, exactitud). Datos exportados a
\texttt{Comparacion\_Final.tex}.}
\label{tab:comp-final}
\begin{tabular}{@{}ccccc@{}}
\toprule
\multirow{2}{*}{Corrida} &
\multicolumn{2}{c}{MLP ($d{=}64$, $M{=}118$)} &
\multicolumn{2}{c}{AlexNet}\\
\cmidrule(lr){2-3}\cmidrule(lr){4-5}
 & Error & Exactitud & Error & Exactitud (\%)\\
\midrule
1  & $0.1208$ & $0.8357$ & $2.16\!\times\!10^{-4}$ & $97.86$\\
2  & $0.0799$ & $0.8929$ & $2.90\!\times\!10^{-4}$ & $97.86$\\
3  & $0.0864$ & $0.8857$ & $5.72\!\times\!10^{-5}$ & $100.00$\\
4  & $0.0879$ & $0.8786$ & $1.25\!\times\!10^{-4}$ & $99.29$\\
5  & $0.1549$ & $0.7786$ & $1.33\!\times\!10^{-4}$ & $98.57$\\
6  & $0.0721$ & $0.8857$ & $6.28\!\times\!10^{-6}$ & $100.00$\\
7  & $0.1198$ & $0.8214$ & $5.27\!\times\!10^{-6}$ & $100.00$\\
8  & $0.1214$ & $0.8286$ & $4.31\!\times\!10^{-9}$ & $100.00$\\
9  & $0.1737$ & $0.7429$ & $2.22\!\times\!10^{-7}$ & $100.00$\\
10 & $0.0996$ & $0.8786$ & $3.10\!\times\!10^{-7}$ & $100.00$\\
\midrule
\textbf{Media} & --- & $\mathbf{0.8429}$ & --- & $\mathbf{99.36}$\\
\textbf{Desv.\ est.} & --- & $0.0486$ & --- & $0.87$\\
\bottomrule
\end{tabular}
\end{table}

La \Cref{fig:comp-curvas} muestra ejemplos de las curvas de entrenamiento de
ambos modelos.

\begin{figure}[htbp]
  \centering
  \begin{subfigure}{0.48\linewidth}
    \centering
    \IfFileExists{figuras/final_mlp_curva.png}
      {\includegraphics[width=\linewidth]{figuras/final_mlp_curva.png}}
      {\fbox{\parbox[c][3.2cm][c]{0.9\linewidth}{\centering\itshape
        \texttt{figuras/final\_mlp\_curva.png}}}}
    \caption{Entrenamiento del MLP.}
  \end{subfigure}\hfill
  \begin{subfigure}{0.48\linewidth}
    \centering
    \IfFileExists{figuras/final_alex_curva.png}
      {\includegraphics[width=\linewidth]{figuras/final_alex_curva.png}}
      {\fbox{\parbox[c][3.2cm][c]{0.9\linewidth}{\centering\itshape
        \texttt{figuras/final\_alex\_curva.png}}}}
    \caption{Entrenamiento de AlexNet.}
  \end{subfigure}
  \caption{Curvas de error de entrenamiento de ambas arquitecturas sobre
    ACRIMA (\texttt{run\_many}).}
  \label{fig:comp-curvas}
\end{figure}

\subsection{Análisis de resultados}

La \Cref{tab:comp-final} evidencia una ventaja sustancial de la arquitectura
convolucional: AlexNet alcanza una exactitud media del $\mathbf{99.36\%}$
(con $7$ de $10$ corridas en $100\%$) frente al $\mathbf{84.29\%}$ del MLP
completamente conectado. Esta diferencia es consistente con la teoría: las capas
convolucionales explotan la estructura espacial de las imágenes mediante
\emph{pesos compartidos} e \emph{invarianza a traslaciones}, mientras que el MLP
trata cada píxel como una característica independiente tras el aplanado,
descartando la vecindad espacial e incrementando enormemente el número de
parámetros. A ello se suma la \emph{transferencia de aprendizaje}: AlexNet parte
de representaciones visuales ya aprendidas en ImageNet, lo que es decisivo ante
la escasez de muestras ($\sim 564$ de entrenamiento). El menor error y la menor
desviación estándar de AlexNet ($0.87$ frente a $4.86$ puntos del MLP) confirman
además una mayor estabilidad entre particiones.

%==============================================================================
\section{Conclusiones}\label{sec:conclusiones}
%==============================================================================

\begin{itemize}[leftmargin=1.4em]
  \item Se implementó, a partir de sus operaciones fundamentales, un perceptrón
        multicapa de tres capas con representación matricial de los pesos,
        entrenado con descenso del gradiente con \emph{momentum}. Las
        modificaciones respecto del \emph{backpropagation} clásico —acumulador
        de velocidad, ponderación convexa $(1-\gamma)$ y gradiente promediado—
        quedan documentadas en \cref{eq:upd_wo,eq:upd_ws,eq:upd_bias}.
  \item La convergencia a error nulo sobre XOR (\Cref{sec:parte2}) valida de
        forma integral la corrección de \emph{forward}, retropropagación y
        actualización, y confirma la capacidad de la red para resolver problemas
        no separables linealmente.
  \item El estudio en $\mathbb{R}^2$ (\Cref{sec:parte3}) mostró, en línea con la
        literatura~\cite{cybenko1989,hornik1989,geman1992}, que $M$ controla la
        capacidad del modelo y que $\alpha$ debe reducirse al aumentar $M$ (y al
        aumentar $\gamma$); con $\alpha=0.01$ las cuatro configuraciones
        convergen a errores próximos, evidenciando una tasa de aprendizaje
        conservadora para el presupuesto de iteraciones empleado.
  \item Sobre ACRIMA (\Cref{sec:parte4}), la arquitectura convolucional AlexNet
        ($99.36\%$ de exactitud media) superó claramente al MLP completamente
        conectado ($84.29\%$), confirmando la idoneidad de las representaciones
        convolucionales y la transferencia de aprendizaje para datos de imagen
        con pocas muestras.
\end{itemize}

\bibliographystyle{IEEEtran}
\bibliography{referencias}

\end{document}
